\section{Coordinates of a product resolution}
\label{sec:orthogonal_resolution_product_coordinates}

Let $I,J,N$ be finite sets. Suppose that
$(P_i)_{i\in I}$ and $(Q_j)_{j\in J}$ are mutually orthogonal
idempotent matrices in $\C^{N\times N}$, each with sum $\Id$.
Assume that $E_{i,j}=P_iQ_j$ also forms an orthogonal resolution of
$\Id$. This holds, in particular, when every $P_i$ commutes with every
$Q_j$. Write
\begin{align*}
  \Phi_P(f)&=\sum_{i\in I}f(i)P_i,&
  \Phi_Q(g)&=\sum_{j\in J}g(j)Q_j,&
  \Phi_E(h)&=\sum_{i\in I}\sum_{j\in J}h(i,j)P_iQ_j.
\end{align*}
The following elementary identities underlie the joint central-label
calculations in Lemma~6.1 and the exponential moment calculations in
Section~7 of~\cite{OpenAI2026AreaLaw}.

\begin{theorem}[Functions of the first coordinate]
  \label{thm:orthogonal_resolution_product_first}
  \lean{Matrix.IsOrthogonalResolution.prod_hom_fst}
  \leanok
  For every $f:I\to\C$,
  \begin{align*}
    \Phi_E\bigl((i,j)\mapsto f(i)\bigr)=\Phi_P(f).
  \end{align*}
\end{theorem}

\begin{proof}
  \leanok
  Summing first over the second coordinate gives
  \begin{align*}
    \sum_i\sum_j f(i)P_iQ_j
      &=\sum_i f(i)P_i\Bigl(\sum_jQ_j\Bigr)
       =\sum_i f(i)P_i.
  \end{align*}
\end{proof}

\begin{theorem}[Functions of the second coordinate]
  \label{thm:orthogonal_resolution_product_second}
  \lean{Matrix.IsOrthogonalResolution.prod_hom_snd}
  \leanok
  For every $g:J\to\C$,
  \begin{align*}
    \Phi_E\bigl((i,j)\mapsto g(j)\bigr)=\Phi_Q(g).
  \end{align*}
\end{theorem}

\begin{proof}
  \leanok
  Interchanging the two finite sums gives
  \begin{align*}
    \sum_i\sum_j g(j)P_iQ_j
      &=\sum_j g(j)\Bigl(\sum_iP_i\Bigr)Q_j
       =\sum_j g(j)Q_j.
  \end{align*}
\end{proof}
